\largerpage
\scp{\tsc{Eastern Islands}}{11-03}
The \tsc{Eastern Islands} node contains \ili{Bati},
\ili{Geser}-\ili{Gorom} of eastern Seram
and adjacent islands, \ili{Watubela}
and \ili{Kesui} of the \ili{Watubela} Islands,
as well as \ili{Kowiai} of Adi Island
and nearby coastal areas of the Bomberai Peninsula
and Bird’s Neck of Papua (see \frf{08-01}).\footnote{
		\ili{As} discussed
		in footnote \ref{fn:SchapperZobel2024}, \citet{SchapperZobel2024b} propose that \ili{Irarutu} (placed
		in \tsc{Tanimbar-Bomberai}, \srf{11-05-02-02}) is most closely related to \ili{Kowiai} -- thus also
		within \tsc{Eastern Islands}.}
The main defining change of \tsc{Eastern Islands} is the lexically
specific split of *j.
Five words have *j > **s in all identified reflexes
and three have *j > **l.
These are shown in \trf{11-13}.
%alongside reflexes in other members of STB for contrast.
Note that the nearby SHWNG subgroup has similar *s/*j > \it{s} \citep[52f]{Kamholz2014}.

\begin{table}\tcp{\tsc{Eastern Islands} *j}{11-13}
\begin{adjustbox}{width=\textwidth}
	\begin{threeparttable}\st{2pt}
	\begin{tabular}{llllll|lll}\lsptoprule
*gloss	&	how m.	&	rice	&	dry	&	nose	&	name	&	sun, day	&	gall	&	ySi\\
PMP	&	*pi\tbr{j}a	&	*pa\tbr{j}ay	&	*ma\tbr{j}a	&	*i\tbr{j}uŋ\su{‡}	&	*ŋa\tbr{j}an	&	*qalə\tbr{j}aw\su{Δ}	&	*qapə\tbr{j}u	&	*hua\tbr{j}i\\	\midrule
\ili{Bati}	&		&	‹faa\tbr{s}í›	&		&	‹\tbr{s}onina›	&		&	‹wo\tbr{l}eh›	&		&	\\
\ili{Geser}	&	{\hr}fi\tbr{s}	&	{\hr}fa\tbr{s}a	&	{\hr}kirima\tbr{s}	&	{\hr}\tbr{s}o	&	{\hr}ŋa\tbr{s}an	&	{\hr}o\tbr{l}a	&	{\hr}fo\tbr{l}i	&	{\hr}a\tbr{l}i\\
\ili{Gorom}	&	{\hr}hi\tbr{s}	&	{\hr}ha\tbr{s}a	&	\cg{(deʔar)}	&	{\hr}i\tbr{s}o, \tbr{s}ue	&	{\hr}ŋa\tbr{s}an	&	{\hr}o\tbr{l}a	&		&	{\hr}a\tbr{l}i\\
\ili{Watubela}	&	{\hr}fi\tbr{h}i	&		&	-ma\tbr{h}a	&	{\hr}\tbr{h}u	&	{\hr}ŋa\tbr{h}an	&		&	{\hr}fe\tbr{l}i-n	&	\\
\ili{Kesui}	&		&	‹fá\tbr{h}a›	&		&		&		&	‹o\tbr{l}ēr›	&		&	\\
\ili{Kowiai}	&	{\hr}ɸi\tbr{s}	&	{\hr}ɸa\tbr{s}a	&	\cg{(namarawar)}\su{†}	&	{\hr}\tbr{s}iʔai	&	{\hr}ne\tbr{s}a	&	{\hr}o\tbr{r}a	&		&	\cg{(ɸuta)}\\	\midrule
\ili{Bobot}	&	{\hr}fila	&		&	{\hr}ma-mala	&	-ilin	&	-nalan	&	{\hr}lia	&	{\hr}foli	&	{\hr}wali-n\\
\ili{Kur}	&		&	{\hr}pasa	&		&		&	{\hr}ŋaa-m	&	{\hr}loŋ, ‹lè›	&		&	\\
\ili{Yamdena}	&	{\hr}fira	&	{\hr}fas-e	&	\cg{(-ŋretw)}	&	{\hr}iri-ŋw	&	{\hr}ŋara-ny	&	{\hr}ler-e	&	{\hr}feru	&	{\hr}wai-ny\\
		\lspbottomrule
	\end{tabular}
		\begin{tablenotes}\tnsize
			\item[†]	{\ili{Kowiai} \it{namarawar} is probably not a reflex of *maja.
								\citet[90]{SmitsVoorhoeve1992b} give the following cognates:
								‹nama rawar›, ‹nama rawa›, ‹namawar›, and ‹namaiawa›.}
			\item[‡]	{If the words given here are indeed reflexes of *ijuŋ,
								most show irregular loss of the initial syllable.
								The source of the final two syllables in \ili{Bati} \it{sonina} is unknown.
								\ili{Kowiai} \it{siʔai} ‘nose’ may be from a different etymon, given the vowels.
								\ili{Gorom} \it{sue} is from \citet{Visser2019a}.
								\ili{Gorom} \it{iso} ‘nose’ is from \citet[23]{Kakerissa1986},
								which is part of the grammar sketch.
								On page 53, in the word list,
								it is given as ‹isoo›.}
			\item[Δ]	{Reflexes of *qaləjaw in \tsc{Eastern Islands} have irregular *l > {\0}.}
		\end{tablenotes}
	\end{threeparttable}
\end{adjustbox}
\end{table}

Additional examples of *j in \tsc{Eastern Islands} include the following,
which also show the split of *j > **s, **l are the following:

\begin{itemize}
	\item *bu\tbr{j}əq > \ili{Gorom} \it{wu\tbr{s}a}, \it{hu\tbr{s}a} ‘foam’
	\item	*pusə\tbr{j} > \ili{Geser} \it{fusa\tbr{l}}
				\ili{Gorom} \it{usa\tbr{l}}, \ili{Kowiai} \it{ɸusa\tbr{r}} ‘navel’
	\item	*bəkəla\tbr{j} > **bəla\tbr{j} > \ili{Kowiai}
				\it{na-ɸora\tbr{s}} ‘spread out (mat)’
	\item	*waRə\tbr{j} > \ili{Kowiai} \it{wara\tbr{s}} ‘rope,
				string’\footnote{
						Another word with *j is *bukij > \ili{Kowiai}
						\it{ɸuʔ} ‘mountain’ which has lost the final syllable.}
\end{itemize}

\tsc{Eastern Islands} languages also share *R > **r word initially
and medially, though depending on the phonetic interpretation
of *R, this may be a retention.
Word finally *R > **r, {\0}.
\ili{Bati} has subsequent initial
and medial **r > \it{l}
and \ili{Watubela} has **r > \it{l} in all positions.
Examples of *R > **r in \tsc{Eastern Islands} are shown in \trf{11-14}.

\begin{table}\tcp{\tsc{Eastern Islands} *R > **r}{11-14}
\begin{adjustbox}{width=\textwidth}
	\begin{threeparttable}\st{2pt}
	\begin{tabular}{lllllllll}\lsptoprule
local gl.	&	house	&	stingray	&	salt	&	ladder	&	bone	&	blood	&	coconut	&	water\\
PMP	&	*\tbr{R}umaq	&	*pa\tbr{R}ih	&	*qasi\tbr{R}a	&	*ha\tbr{R}əzan	&	*du\tbr{R}i	&	*da\tbr{R}aq	&	*niu\tbr{R}	&	*wahiyə\tbr{R}\\	\midrule
\ili{Bati}	&	‹\tbr{l}úme›	&		&	‹si\tbr{l}e›	&		&	‹lu\tbr{l}ul›	&	‹la\tbr{l}ai›	&	‹niū\tbr{l}a›	&	‹a\tbr{rr}›\\
\ili{Geser}	&	{\hr}\tbr{r}uma	&	{\hr}fa\tbr{r}i	&	{\hr}si\tbr{r}a	&	{\hr}\tbr{r}oran	&	{\hr}ru\tbr{r}i	&	{\hr}ra\tbr{r}a	&	{\hr}niu, niu\tbr{r}	&	{\hr}a\tbr{r}\\
\ili{Gorom}	&	{\hr}\tbr{r}uma	&	{\hr}ha\tbr{r}i	&	{\hr}si\tbr{r}a	&	{\hr}\tbr{r}oran	&	{\hr}ru\tbr{r}i	&	{\hr}ra\tbr{r}a	&	{\hr}niu	&	{\hr}a\tbr{r}\\
\ili{Watubela}	&	{\hr}\tbr{l}umak	&	{\hr}fa\tbr{l}i	&		&	{\hr}\tbr{l}alon	&	{\hr}lu\tbr{l}i	&	{\hr}la\tbr{l}ak	&		&	{\hr}a\tbr{l}\\
\ili{Kesui}	&	‹o\tbr{r}úma›	&		&	‹si\tbr{r}a›	&		&	‹lú\tbr{r}u›\su{†}	&	‹la\tbr{r}ah›\su{†}	&	\cg{(‹dar›)}	&	‹a\tbr{rr}›\\
\ili{Kowiai}	&	\cg{(san)}	&	{\hr}ɸa\tbr{r}	&	{\hr}si\tbr{r}a	&	{\hr}\tbr{r}oran	&	{\hr}ru\tbr{r}	&	{\hr}ra\tbr{r}a	&	{\hr}niu, niu\tbr{r}	&	{\hr}wala\tbr{r}\\
		\lspbottomrule
	\end{tabular}
		\begin{tablenotes}\tnsize
			\item[†] {Initial *d > \it{l} in \ili{Kesui} \it{luru} ‘bone’
							and ‹larah› ‘blood’ is probably
							due to regressive dissimilation.}
		\end{tablenotes}
	\end{threeparttable}
\end{adjustbox}
\end{table}

\begin{table}\tcp{\tsc{Eastern Islands} *l}{11-15}
\begin{adjustbox}{width=\textwidth}
	\st{2pt}\begin{tabular}{lllllllll}\lsptoprule
local gl.	&	five	&	sky	&	run	&	egg	&	three	&	eight	&	moon	&	path\\
PMP	&	*\tbr{l}ima	&		*\tbr{l}aŋit	&	*pa-\tbr{l}aRiw	&	*qatə\tbr{l}uR	&	*tə\tbr{l}u	&	*wa\tbr{l}u	&	*bu\tbr{l}an	&	*za\tbr{l}an\\	\midrule
\ili{Bati}	&	‹\tbr{l}im›	&		&		&	‹to\tbr{l}or›	&	‹to\tbr{l}o›	&	‹a\tbr{l}u›	&	‹wúan›	&	‹lāān›\\
\ili{Geser}	&	{\hr}\tbr{l}im	&	{\hr}\tbr{l}aŋit	&		&	{\hr}to\tbr{l}u	&	{\hr}to\tbr{l}u	&	{\hr}a\tbr{l}u	&	{\hr}u\tbr{l}an	&	{\hr}la\tbr{l}an\\
\ili{Gorom}	&	{\hr}\tbr{l}im, lium	&	{\hr}\tbr{l}aŋit	&	{\hr}ha\tbr{l}aru	&	{\hr}to\tbr{l}ur	&	{\hr}to\tbr{l}u	&	{\hr}a\tbr{l}u	&	{\hr}u\tbr{l}an	&	{\hr}la\tbr{l}an\\
\ili{Watubela}	&	{\hr}\tbr{l}ima	&	{\hr}\tbr{l}aŋit	&		&	{\hr}kat\tbr{l}u	&	{\hr}to\tbr{l}u	&		&	{\hr}u\tbr{l}an	&	\\
\ili{Kesui}	&	‹\tbr{r}ima›	&		&		&	‹atu\tbr{l}ú›	&	‹to\tbr{l}u›	&	‹a\tbr{ll}u›	&	‹wú\tbr{l}an›	&	‹\tbr{l}aran›\\
\ili{Kowiai}	&	{\hr}\tbr{r}im	&	{\hr}\tbr{r}aŋgit	&	{\hr}na-fa\tbr{r}ar	&	{\hr}to\tbr{r}on	&	{\hr}to\tbr{r}	&	\cg{(rim tor)}	&	{\hr}ɸu\tbr{r}an	&	{\hr}ra\tbr{r}an\\
		\lspbottomrule
	\end{tabular}
\end{adjustbox}
\end{table}

Two words display *R > {\0} in some reflexes.
Firstly, *baqəRu ‘new’ has *R > {\0} in all identified cases:
\ili{Geser} \it{wou{\tl}wou}; \ili{Gorom} \it{ou{\tl}wou};
\ili{Watubela} \it{oku{\tl}oku}; and \ili{Kowiai} \it{ɸowa}.
Secondly, *maRuqanay ‘male’ has *R > **r in \ili{Geser}-\ili{Gorom} \it{urana}
‘man’ and \ili{Bati} ‹bulana› ‘husband’, but has *R > {\0} in \ili{Watubela} \it{maŋkana}.
\ili{Kowiai} appears to have variant forms.
\citet[137,188]{WalkerWalker1991} have
\it{muana} ‘male’ while \citet[132]{SmitsVoorhoeve1992b} have \it{muruana}.

PMP *l = **l in Eastern Islands with subsequent changes in some daughters.
\ili{Kowiai} has subsequent **l > \it{r,}
and \ili{Bati} has subsequent **l > \it{l,} {\0} /V{\gap}V.
Examples of PMP *l in \tsc{Eastern Islands} are given in \trf{11-15}.
There is also one irregular case of *l > \it{r} in \ili{Kesui}.

The reflexes of PMP *d,
*z, *R, *l, and *j in \tsc{Eastern Islands} are summarised in \trf{11-16}.
A comment on the \ili{Bati} reflexes is in order.
Word medially \ili{Bati} has *z/*l > \it{l},
{\0}, but there are no instances of medial *R > {\0} in the available
data, this suggests that \ili{Bati} *R > \it{l} occurred after cases of *z/*l
> **l > {\0}.
(There are no unambiguous cases of medial *d in the available
\ili{Bati} data to know if it also is lost medially.\footnote{
		There is one possible case of medial *d in the available data: \it{gurun} ‘prawn’.
		While this may be connected with *qudaŋ,
		\ili{Kesui} and \ili{Teor} \it{gurun} ‘prawn’ suggest an irregular local
		development and/or contact.})
Given these complications,
we cannot propose that *d/*z (Proto-STB **d)
and *R merged in Proto-Eastern Islands.
Furthermore, neither of these merged with *l or *j.
These patterns are distinct from those in most \tsc{\ili{Ambon}-Seram}
languages which either show full merger of *d/*z/*R/*l/*j or
only keep *d/*z distinct from *l/*R/*j.\footnote{
		The only \tsc{\ili{Ambon}-Seram}
		languages which do not merge *R with *l/*j
		is Boano in the extreme
		west of \tsc{\ili{Ambon}-Seram} (\srf{10-03-01}),
		as well as \tsc{Seti},
		\ili{Benggoi}, and \ili{Hoti} all of which
		are in the extreme \ili{east} (\srf{10-03-05})
		and thus closer to STB.} 

\begin{table}
	\centering\tcp{\tsc{Eastern Islands} *d, *z, *R, *l, *j}{11-16}
	\st{6pt}\begin{tabular}{llllll@{ }l}\lsptoprule
PMP	&	{\hr}*d	&	{\hr}*z	&	{\hr}*R	&	{\hr}*l	&	{\hr}*j	&	\\
Proto-Eastern Islands	&	**d	&	**d	&	**r	&	**l	&	**s,	&	**l\\	\midrule
\ili{Bati}	&	{\hr\hr}l, ({\0}?)	&	{\hr\hr}l, {\0}	&	{\hr\hr}l	&	{\hr\hr}l, {\0}	&	{\hr\hr}s,	&	{\hr\hr}l\\
\ili{Geser}-\ili{Gorom}	&	{\hr\hr}r	&	{\hr\hr}r	&	{\hr\hr}r	&	{\hr\hr}l	&	{\hr\hr}s,	&	{\hr\hr}l\\
\ili{Watubela}	&	{\hr\hr}l	&	{\hr\hr}l	&	{\hr\hr}l	&	{\hr\hr}l	&	{\hr\hr}h,	&	{\hr\hr}l\\
\ili{Kesui}	&	{\hr\hr}r	&	{\hr\hr}r	&	{\hr\hr}r	&	{\hr\hr}l	&	{\hr\hr}h,	&	{\hr\hr}l\\
\ili{Kowiai}	&	{\hr\hr}r	&	{\hr\hr}r	&	{\hr\hr}r	&	{\hr\hr}r	&	{\hr\hr}s,	&	{\hr\hr}r\\
		\lspbottomrule
	\end{tabular}
\end{table}

The reflexes of *b provide evidence for a node of \tsc{Eastern
Islands} which includes all languages except \ili{Kowiai}.
Before vowels other than \it{u} most cases of *b > \it{w} in
\ili{Bati}, \ili{Geser}-\ili{Gorom}, \ili{Watubela}, and \ili{Kesui}.
\ili{Kowiai} has *b > \it{ɸ} thus merging *p/*b > \it{ɸ}.
Examples are given in \trf{11-17}.
The words for ‘pig’ are irregular in retaining initial *b = \it{b} in all languages.
Additional examples of *b > \it{w} include *buqaya > \ili{Geser}-\ili{Gorom}
\it{wai} ‘crocodile’ and *bujəq > \ili{Gorom} \it{wusa,
husa} ‘foam’.

\begin{table}\tcp{\tsc{Eastern Islands} *b}{11-17}
%\begin{adjustbox}{width=\textwidth}
	\st{2pt}\begin{tabular}{llllllll}\lsptoprule
local gl.	&	stone	&	fly (n.)	&	woman	&	body	&	fruit	&	below	&	pig\\
PMP	&	*\tbr{b}atu	&	*\tbr{b}əRŋaw	&	*(ma)-\tbr{b}inahi	&	*\tbr{b}ataŋ	&	*\tbr{b}uaq	&	*\tbr{b}a\tbr{b}aq	&	*\tbr{b}abuy\\	\midrule
\ili{Bati}	&		&	\cg{(‹langar›)}	&	‹\tbr{b}inei›	&	\cg{(‹rial›)}	&	‹\tbr{w}oya›	&		&	‹\tbr{b}óia›\\
\ili{Geser}	&	{\hr}\tbr{w}atu	&		&	{\hr}wa\tbr{w}ina	&	{\hr}\tbr{w}atan	&	{\hr}\tbr{w}oi	&	{\hr}\tbr{w}a\tbr{w}a	&	{\hr}\tbr{b}oi\\
\ili{Gorom}	&	{\hr}\tbr{w}atu	&	{\hr}\tbr{w}eŋ	&	{\hr}wa\tbr{w}ina	&	{\hr}\tbr{w}atan	&	{\hr}\tbr{w}oiʔ	&	{\hr}\tbr{w}a\tbr{w}a, wa	&	{\hr}\tbr{b}oi\\
\ili{Watubela}	&	{\hr}\tbr{w}atu	&		&	\cg{(hilala)}	&		&	{\hr}okeʔ	&		&	\\
\ili{Kesui}	&		&	‹\tbr{w}éger›	&	\cg{(‹felelára›)}	&	‹\tbr{w}atan›	&	‹\tbr{w}oi›	&		&	‹\tbr{b}oōr›\\
\ili{Kowiai}	&	\cg{(lagera)}	&	{\hr}\tbr{ɸ}eŋ	&	{\hr}me\tbr{ɸ}ina	&	\cg{(ses)}	&	{\hr}aɸa\tbr{ɸ}oi	&	{\hr}ti\tbr{ɸ}aʔa	&	{\hr}\tbr{b}oi\\
		\lspbottomrule
	\end{tabular}
%\end{adjustbox}
\end{table}

Before a high back vowel \it{u} \ili{Kowiai} has *b > \it{ɸ},
while other Eastern Islands languages have *b > {\0},
\it{w} with zero being the most common reflex.
Examples are given in \trf{11-18}.

\begin{table}\tcp{\tsc{Eastern Islands} *bu}{11-18}
%\begin{adjustbox}{width=\textwidth}
	\begin{threeparttable}\st{2.5pt}
	\begin{tabular}{llllllll}\lsptoprule
local gl.	&	fish trap	&	bow	&	moon	&	ten	&	hair	&	sugarcane	&	pig\\
PMP	&	*\tbr{b}u\tbr{b}u	&	*\tbr{b}usuR	&	*\tbr{b}ulan	&	**\tbr{b}utu-sa\su{†}	&	*\tbr{b}uhək	&	*tə\tbr{b}uh	&	*ba\tbr{b}uy\\	\midrule
\ili{Bati}	&		&	‹usulah›	&	‹\tbr{w}úan›	&	‹ocha›	&	‹uka›	&		&	‹bóia›\\
\ili{Geser}	&	{\hr}uu	&	{\hr}usu	&	{\hr}ulan	&	{\hr}uʧa	&	{\hr}uk	&	\cg{(tanu)}	&	{\hr}boi\\
\ili{Gorom}	&	{\hr}au	&	{\hr}\tbr{w}usur	&	{\hr}ulan	&	{\hr}uʧaʔ	&	{\hr}uk, uʔ	&	{\hr}tou	&	{\hr}boi\\
\ili{Watubela}	&	{\hr}uu	&	{\hr}uhul	&	{\hr}ulan	&		&	{\hr}uka	&		&	\\
\ili{Kesui}	&		&	\cg{(‹lóburr›)}	&	‹\tbr{w}úlan›	&	‹sow›	&	‹u′a›	&		&	‹boōr›\\
\ili{Kowiai}	&	{\hr}\tbr{ɸ}u\tbr{ɸ}	&	\cg{(laraʔai)}	&	{\hr}\tbr{ɸ}uran	&	{\hr}\tbr{ɸ}uʧa	&	\cg{(ɸur)}	&	{\hr}to\tbr{ɸ}	&	{\hr}boi\\
		\lspbottomrule
	\end{tabular}
		\begin{tablenotes}\tnsize
			\item[†] {Reflexes of ‘ten’, except perhaps \ili{Kesui},
								are from **butu + *sa ‘one’ > **butsa, with /ts/ → [ʧ].}
		\end{tablenotes}
	\end{threeparttable}
%\end{adjustbox}
\end{table}

A conservative feature of \tsc{Eastern Islands} languages is
the preservation of both the nasal
and a reflex of the plosive in initial *ma-p.
Available examples are: *ma-panas ‘hot’ > \ili{Bati} ‹mofánas›,
\ili{Geser} \it{mfanas}, \ili{Gorom} \it{mahanas},
and *ma-putiq ‘white’ > \ili{Bati} ‹maphutu›,
\ili{Geser} \it{mafuti/futi}, \ili{Gorom} \it{mahuti},
\ili{Watubela} \it{mfutik.}
